% TEMPLATE-MILC-v3.tex - plantilla maestra para todas las guias MILC v3
% Ver scripts/generadores/build-guia-milc.py para la lista completa de
% claves esperadas. El script valida que no queden placeholders sin
% reemplazar antes de escribir el archivo final.

\documentclass[11pt,letterpaper]{article}
\usepackage[margin=1.75cm]{geometry}
\usepackage[table]{xcolor}
\usepackage{fontspec}
\usepackage[spanish]{babel}
\usepackage{tikz}
\usepackage[most]{tcolorbox}
\usepackage{tabularx}
\usepackage{array}
\usepackage{fancyhdr}
\usepackage{titlesec}
\usepackage{ragged2e}
\usepackage{enumitem}
\usepackage{microtype}

\setmainfont{Helvetica Neue}
\setsansfont{Helvetica Neue}
\setlength{\parindent}{0pt}
\setlength{\parskip}{6pt}
\hyphenpenalty=200
\exhyphenpenalty=200
\pretolerance=400
\tolerance=2000
\setlength{\emergencystretch}{1em}
\renewcommand{\arraystretch}{1.25}
\setlist{leftmargin=5mm,itemsep=1mm,topsep=1mm}
\newcolumntype{Y}{>{\RaggedRight\arraybackslash}X}

\definecolor{milcMagenta}{HTML}{E6007E}
\definecolor{milcVino}{HTML}{5A0038}
\definecolor{milcTurquesa}{HTML}{0093A5}
\definecolor{milcVerde}{HTML}{58A923}
\definecolor{milcMostaza}{HTML}{E5B400}
\definecolor{milcOcre}{HTML}{B86600}
\definecolor{milcNegro}{HTML}{241816}
\definecolor{milcGris}{HTML}{F7F7F4}
\definecolor{milcLinea}{HTML}{E8E2DD}
\definecolor{milcGradePrimary}{HTML}{0066FF}
\definecolor{milcGradeDark}{HTML}{003D99}
\definecolor{milcGradeSoft}{HTML}{D6E8FF}
\definecolor{milcGradeAccent}{HTML}{CFE7FF}
\definecolor{milcGradeText}{HTML}{FFFFFF}
\color{milcNegro}

\pagestyle{fancy}
\fancyhf{}
\lhead{\small Tecnología e Informática · Grado 6}
\rhead{\small Guía 21 · MILC v3}
\cfoot{\small\thepage}
\renewcommand{\headrulewidth}{0pt}

\newtcolorbox{titlebox}[1]{
  enhanced,colback=#1,colframe=#1,arc=7mm,boxrule=0pt,
  left=7mm,right=7mm,top=3mm,bottom=3mm,
  before skip=8pt,after skip=8pt,
  fontupper=\sffamily\bfseries\Large\color{white}
}

\newtcolorbox{softbox}[3]{
  enhanced,colback=#2,colframe=#1,borderline west={4pt}{0pt}{#1},
  arc=2mm,boxrule=.4pt,left=4mm,right=4mm,top=3mm,bottom=3mm,
  before skip=6pt,after skip=8pt,title={#3},coltitle=white,
  colbacktitle=#1,fonttitle=\sffamily\bfseries,fontupper=\RaggedRight,
  attach boxed title to top left={xshift=3mm,yshift=-2mm},
  boxed title style={arc=2mm, boxrule=0pt}
}

\newtcolorbox{drawbox}[2]{
  enhanced,colback=white,colframe=#1,arc=4mm,boxrule=1pt,
  left=5mm,right=5mm,top=4mm,bottom=4mm,before skip=8pt,after skip=8pt,
  title={#2},coltitle=white,colbacktitle=#1,fonttitle=\sffamily\bfseries,
  fontupper=\RaggedRight,
  attach boxed title to top left={xshift=4mm,yshift=-2mm},
  boxed title style={arc=3mm, boxrule=0pt}
}

\newtcolorbox{infoband}[1]{
  enhanced,colback=#1!8,colframe=#1!50,arc=2mm,boxrule=.5pt,
  left=4mm,right=4mm,top=2mm,bottom=2mm,before skip=4pt,after skip=4pt,
  fontupper=\small\RaggedRight
}

% Puente narrativo entre secciones (contrato editorial Conéctate).
% Da continuidad: "Acabas de X. Ahora vas a Y porque Z."
% Visualmente: banda izquierda en color del grado, texto en itálica pequeña.
\newtcolorbox{puentebox}{
  enhanced,colback=milcGradeSoft,colframe=milcGradeDark,
  borderline west={3pt}{0pt}{milcGradeDark},
  arc=0mm,boxrule=0pt,left=5mm,right=4mm,top=2.5mm,bottom=2.5mm,
  before skip=4pt,after skip=8pt,
  fontupper=\itshape\small\color{milcNegro}\RaggedRight
}
\newcommand{\puente}[1]{%
  \begin{puentebox}%
    \textcolor{milcGradeDark}{\textbf{\textrightarrow}}\hspace{2mm}#1%
  \end{puentebox}%
}

\newcommand{\linead}[1]{\noindent\color{milcLinea}\rule{#1}{0.5pt}\color{milcNegro}}
\newcommand{\respuesta}[1]{\par\vspace{2mm}\foreach \n in {1,...,#1}{\noindent\color{milcLinea}\rule{\linewidth}{0.5pt}\par\vspace{5mm}}\color{milcNegro}}
\newcommand{\checkbox}{\textcolor{milcGradeDark}{\fbox{\phantom{x}}}\hspace{5pt}}

\newcommand{\phasecard}[4]{
\begin{tcolorbox}[enhanced,colback=#3,colframe=#2,arc=3mm,boxrule=.5pt,height=3.05cm,valign=center,halign=center,left=1.5mm,right=1.5mm,top=2mm,bottom=2mm]
{\sffamily\bfseries\fontsize{22}{24}\selectfont\textcolor{#2}{#1}}\par
{\sffamily\bfseries\small #4}
\end{tcolorbox}
}

\begin{document}
\thispagestyle{empty}
\begin{tikzpicture}[remember picture,overlay]
  \fill[white] (current page.south west) rectangle (current page.north east);
  \path[fill=milcGradePrimary,rounded corners=16mm]
    ([xshift=.55cm,yshift=-.55cm]current page.north west) rectangle
    ([xshift=-.55cm,yshift=3.65cm]current page.south east);

  \node[anchor=north west,text=milcGradeText] at ([xshift=2.0cm,yshift=-1.85cm]current page.north west)
    {\sffamily\bfseries\fontsize{14}{16}\selectfont GRADO};
  \node[anchor=north west,text=milcGradeText,opacity=.82] at ([xshift=2.0cm,yshift=-2.75cm]current page.north west)
    {\sffamily\bfseries\fontsize{9}{11}\selectfont SERIE GUÍAS · TECNOLOGÍA E INFORMÁTICA};

  \begin{scope}[shift={([xshift=-3.05cm,yshift=-2.55cm]current page.north east)}]
    \fill[white,opacity=.80] (-.62,-.52) rectangle (.62,.52);
    \draw[milcGradeText,line width=1pt] (-.62,-.52) rectangle (.62,.52);
    \fill[milcGradeDark] (-.42,-.38) rectangle (-.22,.06);
    \fill[milcGradeAccent] (-.10,-.38) rectangle (.10,.24);
    \fill[milcGradePrimary!60!black] (.22,-.38) rectangle (.42,.38);
  \end{scope}

  \node[anchor=west,text=milcGradeText] at ([xshift=1.9cm,yshift=-12.85cm]current page.north west)
    {\sffamily\bfseries\fontsize{124}{124}\selectfont 6};
  \draw[milcGradeText,line width=1.7pt]
    ([xshift=4.52cm,yshift=-11.68cm]current page.north west) circle (.13cm);

  \node[anchor=west,text=milcGradeText,text width=8.7cm,align=left] at ([xshift=10.35cm,yshift=-11.6cm]current page.north west)
    {
     {\sffamily\bfseries\fontsize{19}{22}\selectfont Apertura\\periodo 3 ---\\escribir y\\buscar con\\criterio}\\[4mm]
     {\sffamily\bfseries\fontsize{23}{26}\selectfont Guía 21-6-TIC}\\[2mm]
     {\sffamily\fontsize{13}{16}\selectfont Período 3 · Procesador de texto e Internet}\\[5mm]
     {\sffamily\bfseries\fontsize{11}{13}\selectfont Institución Educativa Sor María Juliana}\\[1mm]
     {\sffamily\fontsize{10}{12}\selectfont Autor: Dr. Álvaro Cárdenas Orozco}
    };

  \path[fill=milcGradeSoft,rounded corners=7mm]
    ([xshift=1.35cm,yshift=.85cm]current page.south west) rectangle
    ([xshift=-1.35cm,yshift=4.25cm]current page.south east);
  \draw[milcGradeDark,line width=.65pt,rounded corners=7mm]
    ([xshift=1.35cm,yshift=.85cm]current page.south west) rectangle
    ([xshift=-1.35cm,yshift=4.25cm]current page.south east);
  \node[anchor=center,text=milcNegro,text width=17.0cm,align=center] at ([yshift=2.55cm]current page.south)
    {\sffamily\fontsize{10}{12}\selectfont Metodología MILC · Apertura ancestral · Pensamiento computacional · Triángulo Dussel-Estoicismo-Floridi\\
    \textcolor{milcGradeDark}{\bfseries Producto:} En el cuaderno: \textbf{(a) reflexión personal} sobre tu \emph{relación actual con la escritura digital} (¿usas Word? ¿Docs? ¿escribes en celular o computador? ¿cuánto te demoras escribiendo 1 página?); \textbf{(b) reflexión sobre tu uso de internet} (¿cuántas búsquedas en Google haces al día? ¿confías en lo que lees? ¿revisas fuentes?); \textbf{(c) las 5 preguntas-radar} del periodo 3; \textbf{(d) compromiso firmado} de practicar Word + investigar críticamente durante el periodo.};
\end{tikzpicture}
\mbox{}
\newpage

\section*{Apertura: Apertura periodo 3 --- el oficio de escribir y buscar con criterio}

\begin{softbox}{milcGradeDark}{milcGradeSoft}{Saber ancestral}
\textbf{Doña Mercedes la maestra rural de la vereda La Plata de Cartago enseñaba dos cosas: a escribir bien y a leer críticamente.} En su escuelita del campo, los niños aprendían a \textbf{ordenar las ideas en el cuaderno} antes de escribir (\emph{``primero pensar, después escribir''}, decía). Y aprendían a \textbf{verificar las noticias} que llegaban de fuera (\emph{``no todo lo que dice el periódico es verdad''}, repetía). Ese saber doble --- \emph{escribir bien y leer crítico} --- la convirtió en figura respetada de la vereda durante décadas. Sus exalumnos, ya adultos, sabían escribir cartas formales, redactar peticiones, y desconfiar de promesas falsas. \textbf{Hoy, 50 años después, ese mismo saber doble sigue siendo el más útil}, pero la herramienta cambió: ya no es solo cuaderno y periódico, ahora también es Word, Google Docs, Wikipedia, redes sociales. El periodo 3 te entrena en el mismo oficio de doña Mercedes, con las herramientas de tu tiempo.
\end{softbox}

\begin{softbox}{milcTurquesa}{milcGris}{Saber contemporáneo}
En el \textbf{Periodo 3 (Procesador de texto e Internet)} vas a recorrer 10 sesiones que cubren \textbf{dos grandes mitades}. \textbf{Primera mitad (S2 a S6) --- el oficio de escribir digital:} \emph{(S2) Word y Docs; (S3) Formato básico; (S4) Listas y numeración; (S5) Tablas e imágenes; (S6) Encabezado, pie y números}. \textbf{Segunda mitad (S7 a S9) --- el oficio de buscar críticamente:} \emph{(S7) ¿Qué es internet?; (S8) Buscar bien; (S9) Evaluar fuentes}. \textbf{Cierre integrador (S10):} \emph{Mi primer documento informativo} que reúne todo lo aprendido. El periodo 3 es \textbf{el más práctico del año}: cada sesión produce algo concreto en Word o algo concreto en cuaderno sobre internet. Al final del periodo serás capaz de escribir un documento profesional + investigar un tema con criterio. Esas son habilidades para toda la vida adulta: universidad, trabajo, ciudadanía.
\end{softbox}

\begin{softbox}{milcMostaza}{milcGris}{Pregunta-puente}
\textit{Si tu profe te pidiera entregar \emph{``un informe de 2 páginas sobre el río Cauca, formateado con tabla, imagen, fuentes citadas''}, ¿\textbf{sabrías hacerlo} con confianza? Si la respuesta es \emph{``no del todo''}, este periodo es para ti.}
\end{softbox}

\begin{softbox}{milcVerde}{milcGris}{Saber y saber hacer}
Al terminar la clase: (1) podrás \textbf{identificar} la ruta del periodo 3 (las 2 mitades); (2) sabrás \textbf{explicar} por qué escribir y buscar bien son habilidades adultas; (3) podrás \textbf{aplicar} 5 reglas de uso ético de internet; (4) habrás \textbf{creado} reflexión personal + 5 preguntas + compromiso.
\end{softbox}

\small
\begin{tabularx}{\textwidth}{>{\bfseries}p{3.0cm}Y}
\rowcolor{milcGradePrimary!12}Autor & Dr. Álvaro Cárdenas Orozco\\
DBA & Produce contenidos digitales y valida información de fuentes web (MEN, T\&I 6°)\\
Referentes & Saber ancestral de la maestra rural · Lectura crítica · Soberanía informacional\\
Producto final & En el cuaderno: \textbf{(a) reflexión personal} sobre tu \emph{relación actual con la escritura digital} (¿usas Word? ¿Docs? ¿escribes en celular o computador? ¿cuánto te demoras escribiendo 1 página?); \textbf{(b) reflexión sobre tu uso de internet} (¿cuántas búsquedas en Google haces al día? ¿confías en lo que lees? ¿revisas fuentes?); \textbf{(c) las 5 preguntas-radar} del periodo 3; \textbf{(d) compromiso firmado} de practicar Word + investigar críticamente durante el periodo.\\
\end{tabularx}
\normalsize

\puente{Hoy haces 4 cosas: reflexionas sobre tu escritura y búsqueda actuales, conoces la ruta del P3, firmas el compromiso, y defines tus preguntas-radar.}

\begin{titlebox}{milcGradeDark}
Ruta MILC: así vamos a aprender
\end{titlebox}
\begin{tabularx}{\textwidth}{>{\centering\arraybackslash}X>{\centering\arraybackslash}X>{\centering\arraybackslash}X>{\centering\arraybackslash}X}
\phasecard{1}{milcVerde}{milcGris}{Escuta\\[-1mm]\small Observo, escucho y pregunto.} &
\phasecard{2}{milcTurquesa}{milcGris}{Sistematización\\[-1mm]\small Pensamiento computacional.} &
\phasecard{3}{milcMagenta}{milcGris}{Praxis\\[-1mm]\small Construyo y aplico.} &
\phasecard{4}{milcVino}{milcGris}{Evaluación\\[-1mm]\small Reflexión liberadora.}
\end{tabularx}


\puente{Empezamos con un diagnóstico honesto de cómo escribes y buscas hoy.}

\begin{titlebox}{milcVerde}
Fase 1 · Escuta
\end{titlebox}

\begin{softbox}{milcVerde}{milcGris}{Escena para escuchar}
\textbf{Actividad 1 · IDENTIFICA --- Diagnóstico honesto de tu escritura y búsqueda} (12 min · individual). En el cuaderno responde estas 10 preguntas con UNA frase corta cada una. Sé honesto, esto es para que tú mismo veas dónde estás. \textbf{Sobre escritura:} \emph{(1) ¿Cuándo escribiste por última vez algo más largo de 1 página? ¿En cuaderno o digital?} \emph{(2) ¿Has usado Word? ¿Cuántas veces aproximadamente?} \emph{(3) ¿Conoces Google Docs?} \emph{(4) ¿Cuánto te demoras escribiendo 1 página en computador?} \emph{(5) ¿Sabes lo que es ``alineación justificada''?}. \textbf{Sobre búsqueda:} \emph{(6) ¿Cuántas veces al día buscas algo en Google?} \emph{(7) ¿Lees solo el primer resultado o miras varios?} \emph{(8) ¿Has notado anuncios entre los resultados?} \emph{(9) ¿Confías en Wikipedia para tareas?} \emph{(10) ¿Has verificado alguna vez si una noticia que viste en internet es verdadera?}. Anota tus respuestas sin pena: esto es radar, no examen.
\end{softbox}

\checkbox Respondí las 10 preguntas con honestidad.\\
\checkbox No copié ni le pregunté al compañero.\\
\checkbox Me detuve en las preguntas que me dieron duda.

\begin{infoband}{milcVerde}
\textbf{Lo que va al cuaderno (Actividad 1 · IDENTIFICA).} Título: «Actividad 1 --- Mi diagnóstico de escritura y búsqueda». \textbf{Formato:} 10 preguntas con respuesta corta. \textbf{Extensión:} 10 líneas. \textbf{Sabes que terminaste cuando} tus respuestas reflejan tu realidad actual, no una versión imaginada.
\end{infoband}

\puente{Después conoces la ruta del P3 (10 sesiones en 2 mitades) y las 5 reglas de uso ético.}

\begin{titlebox}{milcTurquesa}
Fase 2 · Sistematización con pensamiento computacional
\end{titlebox}

\begin{softbox}{milcTurquesa}{milcGris}{Cuatro pilares aplicados al tema}
Tu diagnóstico te dice dónde estás hoy. Ahora conoce \textbf{la ruta del periodo 3} para saber a dónde llegas en 10 sesiones.
\end{softbox}

\small
\begin{tabularx}{\textwidth}{>{\bfseries\RaggedRight\arraybackslash}p{3.6cm}Y}
\rowcolor{milcTurquesa!90}\color{white}Pilar & \color{white}\bfseries Aplicación al tema\\
1. Descomposición & \textbf{La ruta del periodo 3 --- las 2 mitades.} \textbf{Primera mitad (S2 a S6) --- Procesador de texto:} \emph{S2: ¿Qué es un procesador de texto?} \emph{S3: Formato básico (negrita, alineación, color).} \emph{S4: Listas y numeración.} \emph{S5: Tablas e imágenes.} \emph{S6: Encabezado, pie y números.} Al terminar esta mitad sabes \emph{escribir un documento profesional con todos los elementos}. \textbf{Segunda mitad (S7 a S9) --- Internet:} \emph{S7: ¿Qué es internet? (cómo viajan los datos).} \emph{S8: Buscar bien (operadores de búsqueda).} \emph{S9: Evaluar fuentes (método CRAAP).} Al terminar esta mitad sabes \emph{buscar información y filtrarla con criterio}. \textbf{Cierre integrador (S10):} \emph{Mi primer documento informativo} que combina ambas mitades.\\
2. Reconocimiento de patrones & \textbf{Por qué este periodo es el más práctico.} En P1 viste \emph{conceptos} (identidad digital, netiqueta). En P2 viste \emph{estructura} (hardware, software). En P3 produces \textbf{productos concretos}: documentos, búsquedas, comparaciones de fuentes. \emph{Cada sesión termina con algo que puedes mostrar}. Las habilidades de P3 las vas a usar el resto del año en otras materias (tareas de Historia, informes de Ciencias, etc.) y toda tu vida adulta (universidad, trabajo). No es exagerado: \textbf{escribir bien + buscar bien son las 2 habilidades más rentables del siglo XXI}.\\
3. Abstracción & \textbf{Las 5 reglas de uso ético de internet (no negociables).} (1) \textbf{Cita siempre tus fuentes:} si tomas información de un sitio, dilo. \emph{Copiar sin citar = plagio = falta grave}. (2) \textbf{Verifica antes de creer:} mínimo 2 fuentes confiables que coincidan. (3) \textbf{No comparta lo que no verifica:} si lo dudas, no lo mandes a otros. (4) \textbf{Respeta los derechos de autor:} no usar imágenes con copyright sin permiso. Usa Pixabay, Unsplash. (5) \textbf{Usa internet para sumar, no para restar:} a tu aprendizaje, a otros, al mundo. Si lo que vas a hacer le quita a alguien (su tranquilidad, su reputación, su tiempo), repíensalo.\\
4. Algoritmo conceptual & \textbf{4 herramientas del oficio en P3.} (1) \textbf{Procesador de texto} (Word o Docs): donde produces los documentos. (2) \textbf{Buscador} (Google, Bing): donde investigas. (3) \textbf{Cuaderno físico}: donde planificas antes de escribir digital (\emph{doña Mercedes decía: primero pensar, después escribir}). (4) \textbf{Fuentes confiables}: Wikipedia (punto de partida), sitios de gobierno (.gov.co), universidades (.edu), enciclopedias en línea. Aprenderás a distinguir todo esto en las 10 sesiones.\\
\end{tabularx}
\normalsize

\begin{drawbox}{milcTurquesa}{Ruta P3 + 5 reglas éticas}
\textbf{Ruta P3 (2 mitades):} \\[1mm]
Mitad 1 (Word): S2 (procesador) → S3 (formato) → S4 (listas) → S5 (tablas/imágenes) → S6 (marco). \\[1mm]
Mitad 2 (Internet): S7 (cómo funciona) → S8 (buscar bien) → S9 (evaluar fuentes). \\[1mm]
Cierre: S10 (documento integrador). \\[1mm]
\textbf{5 reglas éticas:} cita fuentes · verifica · no compartas sin verificar · respeta derechos · usa para sumar.
\end{drawbox}

\begin{softbox}{milcMostaza}{milcGris}{Errores comunes y su corrección}
\textbf{Errores frecuentes al iniciar P3.} (1) \emph{``Yo ya sé Word, este periodo es repetir lo que sé''} --- Falso. Word tiene mucho más de lo que la mayoría usa. Las S2 a S6 te enseñan funciones profesionales. (2) \emph{``Buscar en Google es solo escribir y leer el primer resultado''} --- Falso. Hay técnicas (operadores, filtros, evaluación de fuentes) que multiplican tu eficiencia. (3) \emph{``Copiar de Wikipedia está mal''} --- Discutible. Wikipedia es buen punto de partida, lo malo es \emph{citar Wikipedia} en un trabajo formal. Citas la fuente original que Wikipedia referencia. (4) \emph{``Si está en internet, es cierto''} --- Falso, en 2026 más que nunca. Verificar es habilidad clave.
\end{softbox}

\begin{infoband}{milcTurquesa}
\textbf{Lo que va al cuaderno (Actividad 2 · EVALÚA).} Título: «Actividad 2 --- Ruta del P3 + me comprometo o no con cada regla ética». \textbf{Formato:} bloque A con la ruta del P3 (mitades). Bloque B con las 5 reglas éticas, y al lado de cada una decide \textbf{Sí me comprometo / No estoy seguro} con una razón corta (no palomita ciega: decisión razonada). \textbf{Veredicto (2 líneas):} ¿\emph{cuál} de las 5 reglas te parece la más difícil de cumplir y por qué? Tu honestidad importa más que tu palomita. \textbf{Extensión:} 15 líneas + veredicto. \textbf{Sabes que terminaste cuando} pudiste decir ``No estoy seguro'' en al menos una regla sin sentir vergüenza --- esa es la regla a trabajar este periodo.
\end{infoband}

\puente{Con la ruta clara, escribes tu reflexión personal y firmas el compromiso.}

\begin{titlebox}{milcMagenta}
Fase 3 · Praxis con pensamiento computacional
\end{titlebox}

\begin{softbox}{milcMagenta}{milcGris}{Construyendo el producto con los cuatro pilares}
\textbf{Actividad 3 · CREA --- Reflexión personal + 5 preguntas + compromiso} (25 min · individual). Vas a producir 3 cosas en el cuaderno: una reflexión sobre tu uso actual, 5 preguntas-radar para el periodo, y un compromiso firmado.
\end{softbox}

\small
\begin{tabularx}{\textwidth}{>{\bfseries\RaggedRight\arraybackslash}p{3.6cm}Y}
\rowcolor{milcMagenta!90}\color{white}Pilar & \color{white}\bfseries Aplicación al producto\\
Descomposición & \textbf{Paso 1 --- Reflexión personal (mitad superior de la hoja).} Toma tu diagnóstico de Actividad 1 (las 10 preguntas) y \textbf{convierte las 10 respuestas en un párrafo reflexivo} de 5-6 líneas. Empieza con \emph{``Hoy en escritura digital yo...''} y termina con \emph{``Lo que quiero mejorar este periodo es...''}. No copies las respuestas: refléjalas. Esto te entrena en redactar tu propia voz.\\
Patrones & \textbf{Paso 2 --- Las 5 preguntas-radar del periodo (centro de la hoja).} Subtítulo: \textbf{``Mis 5 preguntas-radar del P3''}. Lista numerada de 5 preguntas que esperas resolver durante el periodo. Sugerencias para inspirarte (no copies estas, inventa tuyas): \emph{``¿Cómo escribo un documento que se vea profesional?''}, \emph{``¿Cómo busco información sin caer en publicidad?''}, \emph{``¿Cómo sé si una página dice verdad?''}, \emph{``¿Cómo cito fuentes correctamente?''}, \emph{``¿Cuántas veces al día debería revisar las fuentes que leo?''}.\\
Abstracción & \textbf{Paso 3 --- Las 5 reglas éticas firmadas (mitad inferior).} Subtítulo: \textbf{``Mi compromiso ético del P3''}. Reproduce las 5 reglas (cita fuentes, verifica, no comparta sin verificar, respeta derechos, usa para sumar). Al lado de cada una, palomita ✓ si te comprometes. Si dudas con alguna, conversa con el profe antes de firmar.\\
Algoritmo de armado & \textbf{Paso 4 --- Compromiso del oficio.} Debajo de las reglas, escribe esta frase en grande: \textbf{``Me comprometo a escribir con mi voz, no con voz prestada. A buscar con criterio, no con prisa.''}. Subráyala. Estos son los 2 pilares éticos del periodo 3.\\
\end{tabularx}
\normalsize

\begin{drawbox}{milcMagenta}{Antes de cerrar la S1 del P3}
\checkbox La reflexión personal es un párrafo de 5-6 líneas escrito con mi voz. \\[1mm]
\checkbox Tengo 5 preguntas-radar reales mías. \\[1mm]
\checkbox Las 5 reglas éticas están firmadas. \\[1mm]
\checkbox El compromiso del oficio está subrayado. \\[1mm]
\checkbox La sesión 1 tiene encabezado y registro de las 3 actividades. \\[1mm]
\checkbox Firmé con nombre completo y fecha.
\end{drawbox}

\begin{softbox}{milcMagenta}{milcGris}{Plantilla de trabajo}
\textbf{Estructura.} \textit{Mitad superior:} reflexión personal (5-6 líneas). \textit{Centro:} 5 preguntas-radar. \textit{Mitad inferior:} 5 reglas éticas firmadas + compromiso del oficio subrayado. \textit{Cuaderno principal:} sesión 1 con encabezado fijo.
\end{softbox}

\begin{infoband}{milcMagenta}
\textbf{Lo que va al cuaderno (Actividad 3 · CREA).} Título: «Actividad 3 --- Mi entrada al periodo 3». \textbf{Formato:} reflexión + 5 preguntas + reglas firmadas + sesión 1 con encabezado. \textbf{Extensión:} 1.5 hojas. \textbf{Sabes que terminaste cuando} tu cuaderno está listo para empezar la S2 con confianza.
\end{infoband}

\puente{El producto es la reflexión + 5 preguntas-radar + compromiso firmado.}

\begin{titlebox}{milcMostaza}
Producto de la sesión
\end{titlebox}
\textbf{Reflexión + 5 preguntas + compromiso ético del periodo 3 · firmado}.
\begin{softbox}{milcMostaza}{milcGris}{Criterios de calidad}
(1) La reflexión personal está escrita con voz propia (5-6 líneas). (2) Las 5 preguntas-radar son personales y específicas. (3) Las 5 reglas éticas tienen firma de compromiso. (4) El compromiso del oficio está subrayado y visible. (5) La sesión 1 tiene encabezado fijo y registro de actividades.
\end{softbox}

\puente{Antes de salir, verifica que tus 5 preguntas son sinceras (no copiadas).}

\begin{titlebox}{milcVino}
Fase 4 · Evaluación liberadora — cinco dimensiones
\end{titlebox}

\begin{softbox}{milcVino}{milcGris}{Sentido de la evaluación}
Evaluar no es solo revisar una nota. Es reconocer qué aprendí, qué cambió en mí, qué emociones manejé, cómo me relaciono con la ciudad y la comunidad, y qué le dejaré a quien viene detrás.
\end{softbox}

\textbf{Mi reflexión liberadora en cinco dimensiones.} Completa cada frase iniciada:

\small
\begin{tabularx}{\textwidth}{>{\bfseries\RaggedRight\arraybackslash}p{3.4cm}Y>{\RaggedRight\arraybackslash}p{4.6cm}}
\rowcolor{milcVino}\color{white}Dimensión & \color{white}\bfseries Mi respuesta (frame starter) & \color{white}\bfseries Algo que descubrí\\
1. Personal & Lo que más cambió en mí fue \linead{6.5cm} & \linead{4.2cm}\\
2. Emocional & La emoción más fuerte fue \linead{2.5cm} y la regulé \linead{3cm} & \linead{4.2cm}\\
3. Ciudadana & En democracia esto importa porque \linead{6cm} & \linead{4.2cm}\\
4. Local (Cartago) & En mi barrio sería útil para \linead{6.5cm} & \linead{4.2cm}\\
5. Intergeneracional & A mi abuela/abuelo le diría \linead{6.5cm} & \linead{4.2cm}\\
\end{tabularx}
\normalsize

\begin{infoband}{milcVino}
\textbf{Para la próxima sustentación.} Vuelve a esta tabla en seis meses. Verás que las respuestas cambian — no porque lo aprendido cambie, sino porque tú cambias. La evaluación liberadora es una conversación contigo a través del tiempo.
\end{infoband}


\puente{Cerramos con 3 voces que te ayudan a entender por qué saber escribir y buscar es ciudadanía moderna.}

\begin{titlebox}{milcGradeDark}
Triángulo de pensamiento · cierre filosófico
\end{titlebox}

\begin{softbox}{milcVino}{milcGris}{Dussel · lente del nosotros}
\textit{````Escribir con tu voz es decir: existo, pienso, tengo algo que aportar. Eso es soberanía.''''}

Mucha gente vive sin escribir nunca su propia voz: copia, repite, replica. \textbf{El que escribe con voz propia se vuelve sujeto, no objeto}. El que solo recibe información sin producir queda dependiente. \emph{Este periodo te entrena en producir, no solo consumir}. Es la diferencia entre ser leído y ser leyente. Doña Mercedes lo sabía: enseñar a escribir es enseñar a ser persona.

\textbf{¿Estoy produciendo voz propia o solo consumiendo voces ajenas?}
\end{softbox}
\linead{\linewidth}\\[1mm]
\linead{\linewidth}

\begin{softbox}{milcMostaza}{milcGris}{Estoicismo · Marco Aurelio (emperador que escribía con disciplina) · cuidado interior}
\textit{````Escribir es pensar despacio. Buscar es elegir despacio. Las dos disciplinas hacen al ciudadano libre.''''}

\textbf{Escribir y buscar son disciplinas lentas} en un mundo rápido. La gente apurada \emph{escribe mal y busca peor}: improvisa, copia, se cree el primer resultado. La gente disciplinada \emph{escribe con cuidado y busca con criterio}: produce documentos profesionales y encuentra información buena. La diferencia se nota al cabo de los años: unos quedan dependientes, otros toman decisiones por sí mismos.

\textbf{¿Me estoy entrenando en escribir y buscar despacio, o sigo el ritmo apurado de los demás?}
\end{softbox}
\linead{\linewidth}\\[1mm]
\linead{\linewidth}

\begin{softbox}{milcTurquesa}{milcGris}{Floridi · lente de la infoesfera}
\textit{````En el siglo XXI, no saber escribir bien y no saber buscar bien es como en el siglo XX no saber leer. Te excluye de decisiones importantes.''''}

\textbf{Tu vida adulta dependerá de escribir y buscar bien}: para postular trabajos, hacer trabajos universitarios, defender derechos, votar con criterio, criar hijos en una era digital. \emph{Las habilidades del P3 son alfabetización del siglo XXI}. No es opcional. Tu yo de 25 años te lo agradecerá. Tu yo de 50 años no podría vivir sin esto.

\textbf{¿Estoy desarrollando las habilidades que mi yo de 25 años va a necesitar para una vida digna?}
\end{softbox}
\linead{\linewidth}\\[1mm]
\linead{\linewidth}

\begin{titlebox}{milcVino}
Compromiso final
\end{titlebox}

\small
\begin{tabularx}{\textwidth}{>{\RaggedRight\arraybackslash}p{3.0cm}YYY}
\rowcolor{milcVino}\color{white}\bfseries Criterio & \color{white}\bfseries Lo logré & \color{white}\bfseries En proceso & \color{white}\bfseries Necesito apoyo\\
Comprensión del tema & Explico ideas centrales con mis palabras. & Reconozco ideas pero falta precisión. & Necesito repasar conceptos base.\\
Pensamiento computacional & Apliqué los 4 pilares al diseñar mi producto. & Apliqué algunos pilares. & Necesito releer la fase 2.\\
Aplicación y producto & Mi producto cumple los criterios y se puede revisar. & Producto funcional con ajustes pendientes. & Aún en construcción.\\
Reflexión 5 dimensiones & Respondí cada dimensión con honestidad. & Respondí algunas con detalle. & Mis respuestas son muy generales.\\
Triángulo de pensamiento & Conecté las 3 voces con mi trabajo real. & Conecté al menos una. & No relaciono las voces con mi proceso.\\
\end{tabularx}
\normalsize

\textbf{Hoy entendí que:}\\[1mm]
\linead{\linewidth}\\[1mm]
\linead{\linewidth}

\textbf{Mi compromiso para la próxima sesión es:}\\[1mm]
\linead{\linewidth}\\[1mm]
\linead{\linewidth}

\begin{softbox}{milcMostaza}{milcGris}{Cita de despedida · Marco Aurelio}
\textit{``El obstáculo en el camino se vuelve el camino. Lo que estorba al camino se vuelve camino.''}\\[1mm]
Cada error de hoy, bien observado, se convierte en pieza del camino que pisarás mañana.
\end{softbox}

\small
\begin{tabularx}{\textwidth}{>{\bfseries}p{3.6cm}Y}
\rowcolor{milcGradeDark!12}Nombre completo: & \linead{10cm}\\
Fecha: & \linead{4cm} \quad Curso/grupo: \linead{4cm}\\
Mi firma: & \linead{9cm}\\
\end{tabularx}
\normalsize

\begin{infoband}{milcGradeDark}
\textbf{Para la próxima sesión.} Trae el cuaderno listo. La próxima clase es S2: \emph{¿Qué es un procesador de texto?}. Vas a crear tu primer documento digital de 3 párrafos sobre un tema que te guste. Recuerda lo que dijo doña Mercedes: \emph{``primero pensar, después escribir''}. Trae una idea de tema para empezar de una vez.
\end{infoband}

\end{document}
